\chapter{Quadratic algebras of order \(9\)}
\label{chap:algebras-cuadraticas-orden-nueve}
\index{quadratic algebra}
\index{finite field \(\mathbb F_9\)}

The identity \(A_W^2=-I_6\), obtained in the preceding section by reducing
the Paley conference matrix, contains more information than a
modular orthogonality relation. It determines a scalar extension of degree two,
turns the six-dimensional ternary space into a three-dimensional space
over the field of nine elements, and permits the ternary Golay code to be
written as the graph of a quadratic unit. This section develops
those consequences and distinguishes them from the characteristic-nine arithmetic
used by the Hensel branch.

\section{Estructura of \texorpdfstring{\(\mathbb F_9\)}{F9} inducida by Paley}

The polynomial \(X^2+1\) has no roots in \(\mathbb F_3\): its values in
\(0,1,2\) are \(1,2,2\). It is therefore irreducible, and the quotient
\[
 \mathbb F_9
 :=\mathbb F_3[\iota]/(\iota^2+1)
\]
is a field of nine elements.

\begin{theorem}[Structure quadratic induced by \(A_W\)]
\label{thm:f9-tres-desde-paley}
Let \(V=\mathbb F_3^6\), written by means of row vectors, and let
\(J(v)=vA_W\). The rule
\[
 (a+b\iota)v:=av+bJ(v),
 \qquad a,b\in\mathbb F_3,
\]
turns \(V\) into a vector space of dimension three over
\(\mathbb F_9\). In the coordinates fixed by the Paley construction,
\[
 \mathbb F_3^6\cong\mathbb F_9^3.
\]
\end{theorem}

\begin{proof}
The equality \(J^2=-I\) makes the evaluation
\(\mathbb F_3[X]\to\operatorname{End}_{\mathbb F_3}(V)\), \(X\mapsto J\),
vanish on the ideal \((X^2+1)\). The action therefore factorizes through
the field \(\mathbb F_9\). Since
\([\,\mathbb F_9:\mathbb F_3\,]=2\) and
\(\dim_{\mathbb F_3}V=6\), one obtains
\(\dim_{\mathbb F_9}V=3\).
\end{proof}

The isomorphism depends on the coordinate system that fixes \(A_W\). The
assertion does not follow from the mere cardinal \(3^6=9^3\): the matrix exhibits
multiplication by all scalars \(a+b\iota\) and therefore provides the
complete linear structure.

\section{The ternary code as an isotropic graph}

The code's generating application is
\[
 c(q)=(q,qA_W).
\]
After identifying \(J\) with multiplication by \(\iota\), its image
admits the writing
\[
 C_W=\{(q,\iota q):q\in\mathbb F_9^3\},
\]
where both components are regarded, by restriction of scalars, as
vectors of \(\mathbb F_3^6\).

\begin{proposition}[Isotropic maximality]
\label{prop:golay-grafo-isotropico}
With respect to the standard bilinear form of \(\mathbb F_3^{12}\), the
subspace \(C_W\) is totally isotropic and maximal among the totally
isotropic subspaces.
\end{proposition}

\begin{proof}
The generator matrix \(G=(I_6\mid A_W)\) satisfies
\[
 GG^{\mathsf T}=I_6+A_WA_W^{\mathsf T}=I_6+2I_6=0.
\]
The generated space is therefore totally isotropic. It has
dimension six, half the dimension of the nondegenerate ambient space,
and is therefore maximal; indeed, it coincides with its orthogonal dual.
\end{proof}

The root of \(-I\) thus appears as a finite incidence transformation.
The analytic complex field is not introduced into the system: a
quadratic structure over \(\mathbb F_3\) is obtained, sufficient to organize the
polarity, the code, and its two eigenspaces after extending
scalars.

\section{Characteristic \(9\) and characteristic \(3\)}

Two sets of \(729\) states enter the construction:
\[
 (\mathbb Z/9\mathbb Z)^3
 \qquad\text{and}\qquad
 \mathbb F_9^3.
\]
They are not isomorphic as additive groups or as rings. The first has
characteristic nine and contains elements of order nine; the second has
characteristic three and its entire additive group has exponent three.

For \(a,b\in\mathbb F_3\), choosing representatives
\(\widetilde a,\widetilde b\in\{0,1,2\}\), let
\[
 \beta(a,b)=\widetilde a+3\widetilde b\pmod 9.
\]
The coordinate map induced by \(\beta\) is a bijection of sets between \(\mathbb F_3^6\) and \((\mathbb Z/9\mathbb Z)^3\), but does not preserve addition. Indeed, if
\[
 \kappa(a,c)
 :=\left\lfloor
 \frac{\widetilde a+\widetilde c}{3}
 \right\rfloor,
\]
then
\[
 \beta(a,b)+\beta(c,d)
 =\beta\bigl(a+c,\ b+d+\kappa(a,c)\bigr),
\]
with the coordinates of the second member reduced modulo three. For example,
\[
 \beta(2,0)+\beta(2,0)=\beta(1,1),
\]
whereas coordinatewise addition in \(\mathbb F_3^2\) would give \((1,0)\). The
second coordinate of \(\beta(1,1)\) preserves the quotient of Euclidean
division; this term is precisely the datum lost under
componentwise reduction.

This distinction types the two readings of the cardinal \(729\). The realization
\((\mathbb Z/9\mathbb Z)^3\) preserves residue and quotient; the realization
\(\mathbb F_9^3\) converts each ternary pair into a scalar of a quadratic
extension. One prolongs arithmetic with division memory; the other
organizes exceptional incidence.

\begin{proposition}[Two typed readings of the same ternary support]
\label{prop:dos-lecturas-tipadas-729}
Once the set-theoretic section used in \(\beta\), the basis, the
polarization, and the operator \(J\) of
\cref{thm:f9-tres-desde-paley} are fixed, there exists the diagram of set bijections
\[
\boxed{
(\mathbb Z/9\mathbb Z)^3
\xleftarrow{\;\beta^{\times3}\;}
\mathbb F_3^6
\xrightarrow{\;\iota_J\;}
\mathbb F_9^3.}
\]
The left arrow transports the nonadic sum twisted by the carry
cocycle \(\kappa\); the right arrow transports the linear structure over
\(\mathbb F_9\) determined by \(J^2=-I\). The two arrows share the
coordinate support \(\mathbb F_3^6\), but do not identify their algebraic
laws.
\end{proposition}

\begin{proof}
The formula
\[
 \beta(a,b)+\beta(c,d)
 =\beta\bigl(a+c,b+d+\kappa(a,c)\bigr)
\]
proves that \(\beta^{\times3}\) is a bijection with twisted addition and is not
an isomorphism from the componentwise sum of
\(\mathbb F_3^6\). By \cref{thm:f9-tres-desde-paley}, the choice of
\(J\) turns that same six-dimensional ternary space into a
three-dimensional space over
\(\mathbb F_3[J]\cong\mathbb F_9\), thereby defining \(\iota_J\).
\end{proof}

The proposition is a synthesis demonstrated with explicit starting structure.
It does not assert a canonical identification between
\((\mathbb Z/9\mathbb Z)^3\) and \(\mathbb F_9^3\): the first contains
additive elements of order nine, whereas the second has additive exponent
three. The cocycle measures precisely the information that would be lost by
confusing both laws.

\section{Classification of Quadratic Algebras}

Over a field of characteristic other than two, a unital quadratic algebra
is classified by the type of its defining polynomial. Over
\(\mathbb F_3\) the three representatives may be chosen as
\[
 \mathcal Q_{\mathrm e}
 =\mathbb F_3[t]/(t^2+1)\cong\mathbb F_9,
\]
\[
 \mathcal Q_{\mathrm n}
 =\mathbb F_3[\varepsilon]/(\varepsilon^2),
 \qquad
 \mathcal Q_{\mathrm s}
 =\mathbb F_3[r]/(r^2-1)
 \cong\mathbb F_3\times\mathbb F_3.
\]

\begin{theorem}[Classification of the quadratic algebras over \(\mathbb F_3\)]
\label{thm:clasificacion-cuadratica-nonadica}
The relations
\[
 J^2=-I,\qquad N^2=0,\qquad R^2=I
\]
define actions of
\(\mathcal Q_{\mathrm e}\), \(\mathcal Q_{\mathrm n}\), and
\(\mathcal Q_{\mathrm s}\), respectively. These three algebras have
nine elements and are mutually nonisomorphic: the first is a field; the
second contains nonzero nilpotents; the third contains nontrivial
idempotents.
\end{theorem}

\begin{proof}
Each relation makes the evaluation of
\(\mathbb F_3[t]\) factor through the corresponding ideal. The irreducibility of
\(t^2+1\) proves that \(\mathcal Q_{\mathrm e}\) is a field. In
\(\mathcal Q_{\mathrm n}\), the class \(\varepsilon\) is nonzero and satisfies
\(\varepsilon^2=0\). In \(\mathcal Q_{\mathrm s}\), since \(2^{-1}=2\) in
\(\mathbb F_3\), the elements
\[
 e_+=2(1+r),\qquad e_-=2(1-r)
\]
are nontrivial orthogonal idempotents and sum to the unit. Being a field,
possessing a nonzero nilpotent, and possessing nontrivial idempotents are
properties invariant under isomorphism.
\end{proof}

Three copies of the two-dimensional regular representation yield
three six-dimensional actions on \(\mathbb F_3\). The relation
\(J^2=-I\) corresponds to Paley polarity; relation \(N^2=0\)
describes the linearization of an extension with quotient; relation
\(R^2=I\) separates two summands through the idempotents \(e_\pm\). This
classification supplies the finite algebraic framework in which the
elliptic, nonreduced, and split structures are distinguished without appeal to a
continuous geometry.

The construction does not replace \(\mathbb C\): the field \(\mathbb F_9\) does not
possess the topology, completeness, or complex analytic structure. The
exact result is different: a square root of \(-I\) and the corresponding
extension of scalars emerge from the incidence matrix already constructed
within the ternary system.

\section{Integral quadratic form and finite and real specializations}
\label{sec:fuente-cuadratica-integral-rev3}

The finite root and the real operator root are coordinated through a common integral presentation. Let
\begin{equation}
 \mathcal Q_{\mathbb Z}
 :=\mathbb Z[u]/(u^2+1).
 \label{eq:fuente-cuadratica-integral-rev3}
\end{equation}
Its change of basis are
\[
 \mathcal Q_{\mathbb Z}\otimes_{\mathbb Z}\mathbb F_3
 \cong\mathbb F_3[u]/(u^2+1)
 \cong\mathbb F_9,
\]
and
\[
 \mathcal Q_{\mathbb Z}\otimes_{\mathbb Z}\mathbb R
 \cong\mathbb R[u]/(u^2+1).
\]

\begin{theorem}[Two realizations of an integral quadratic form]
\label{thm:dos-realizaciones-fuente-cuadratica-rev3}
The finite representation \(u\mapsto A_W\) on
\(\mathbb F_3^6\) and the real representation \(u\mapsto J_{\mathrm e}\) on
\(V_{\mathrm e}\), constructed in the
\cref{thm:tres-generadores-concretos-trit}, realize, after their
respective changes of basis, the same integral relation
\[
 u^2+1=0.
\]
\end{theorem}

\begin{proof}
The equality \(A_W^2+I_6=0\) makes the evaluation
\(\mathbb F_3[u]\to\operatorname{End}_{\mathbb F_3}(\mathbb F_3^6)\) factor through
\((u^2+1)\). In the same way, \(J_{\mathrm e}^2+I=0\) makes
\(\mathbb R[u]\to\operatorname{End}_{\mathbb R}(V_{\mathrm e})\) factor. The two
representations follow from specializations of
\(\mathcal Q_{\mathbb Z}\); no unital homomorphism
\(\mathbb F_9\to\mathbb R\) exists or is used.
\end{proof}

The common source coordinates relations; it does not identify modules of
distinct characteristics. The finite branch preserves incidence; the real branch
incorporates the analytic structure required for the operator exponential.

\section{\(3\) generators concrete of the TRIT regimes}
\label{sec:tres-generadores-concretos-trit}
\index{TRIT!concrete generators}
\index{elliptic generator}
\index{parabolic generator}
\index{hyperbolic generator}

The \cref{chap:trit-algebras} fixes the three abstract quadratic types
\(J^2=-I\), \(J^2=0\), and \(J^2=I\), while the
\cref{thm:clasificacion-cuadratica-nonadica} provides their finite models
on \(\mathbb F_3\). Those objects must not be confused with the
subsequent real realizations. Here three real spaces and three concrete
operators are fixed, each obtained from its declared causal antecedent.

Let
\[
V_{\mathrm e}=A_2\otimes_{\mathbb Z}\mathbb R,
\qquad
V_{\mathrm p}=\mathbb R^2,
\qquad
V_{\mathrm h}=\mathbb R^2.
\]
On \(V_{\mathrm e}\) acts the Coxeter \(C\) of the
\cref{prop:coxeter-a2-primario}. On \(V_{\mathrm p}\) it fixes
\[
N=
\begin{pmatrix}
0&1\\
0&0
\end{pmatrix}.
\]
Finally, the \cref{thm:descenso-necesario-fav} constructs
\[
F_{\rm av}=
\begin{pmatrix}
0&1\\
1&1
\end{pmatrix},
\qquad
A:=F_{\rm av}^{\,2}=
\begin{pmatrix}
1&1\\
1&2
\end{pmatrix}
\]
on \(V_{\mathrm h}\).

\begin{theorem}[Three concrete realizations of the TRIT regimes]
\label{thm:tres-generadores-concretos-trit}
Define
\[
J_{\mathrm e}:=\frac{2C+I}{\sqrt3},
\qquad
J_{\mathrm p}:=N,
\qquad
J_{\mathrm h}:=\frac{2A-3I}{\sqrt5}.
\]
Then
\[
J_{\mathrm e}^{\,2}=-I_{V_{\mathrm e}},
\qquad
J_{\mathrm p}^{\,2}=0,
\qquad
J_{\mathrm h}^{\,2}=I_{V_{\mathrm h}}.
\]
Therefore, for \(t\in\mathbb R\),
\[
\exp(tJ_{\mathrm e})
=\cos t\,I+\sin t\,J_{\mathrm e},
\qquad
\exp(tJ_{\mathrm p})
=I+tJ_{\mathrm p},
\]
\[
\exp(tJ_{\mathrm h})
=\cosh t\,I+\sinh t\,J_{\mathrm h}.
\]

Furthermore, in the elliptic branch,
\[
C=-\frac12I+\frac{\sqrt3}{2}J_{\mathrm e},
\]
and, for \(n\in\mathbb Z\),
\[
C^n=
\cos\!\left(\frac{2\pi n}{3}\right)I
+\sin\!\left(\frac{2\pi n}{3}\right)J_{\mathrm e}.
\]
In the hyperbolic branch, the form
\[
G=
\begin{pmatrix}
2&1\\
1&-2
\end{pmatrix}
\]
has signature \((1,1)\) and satisfies
\[
A^{\mathsf T}GA=G,\qquad
\det A=1,\qquad
\operatorname{tr}A=3.
\]
If
\[
\eta:=\operatorname{arcosh}\frac32=2\log\varphi,
\]
then
\[
A=\cosh(\eta)I+\sinh(\eta)J_{\mathrm h},
\qquad
A^n=\cosh(n\eta)I+\sinh(n\eta)J_{\mathrm h}.
\]
\end{theorem}

\begin{proof}
The polinomio caracteristico of the Coxeter is \(X^2+X+1\); by
Cayley--Hamilton,
\[
C^2+C+I=0.
\]
From here,
\[
(2C+I)^2=4C^2+4C+I=-3I,
\]
and, therefore, \(J_{\mathrm e}^2=-I\). The identity for \(C^n\) follows from
\(C=-\frac12I+\frac{\sqrt3}{2}J_{\mathrm e}\).

Direct multiplication gives \(N^2=0\). Hence, the exponential series
truncates exactly:
\[
\exp(tN)=I+tN.
\]

For the third branch,
\[
A^{\mathsf T}GA=G,\qquad
\det G=-5<0,
\]
so that \(G\) is nondegenerate and has signature \((1,1)\). Likewise,
\[
(2A-3I)^2=5I,
\]
therefore \(J_{\mathrm h}^2=I\). As
\[
\cosh(2\log\varphi)=\frac32,
\qquad
\sinh(2\log\varphi)=\frac{\sqrt5}{2},
\]
we obtain
\[
A=\frac32I+\frac{\sqrt5}{2}J_{\mathrm h}
=\cosh(\eta)I+\sinh(\eta)J_{\mathrm h}.
\]
The formula for \(A^n\) follows from the exponential law. The three general
expressions are obtained, in each case, by separating the even and odd
powers in the exponential series.
\end{proof}

\begin{corollary}[Extended Euler equation by the tritic parameter]
\label{cor:euler-extendida-trit-fibonacci}
For \(\tau\in\{+1,0,-1\}\), let \(J_\tau\) be one of the preceding generators, normalized by
\[
 J_\tau^2=-\tau I.
\]
Define the functions of the quadratic family.
\[
 C_\tau(t)=
 \begin{cases}
  \cos t,&\tau=+1,\\
  1,&\tau=0,\\
  \cosh t,&\tau=-1,
 \end{cases}
 \qquad
 S_\tau(t)=
 \begin{cases}
  \sin t,&\tau=+1,\\
  t,&\tau=0,\\
  \sinh t,&\tau=-1.
 \end{cases}
\]
Then the three branches satisfy a single exponential identity:
\begin{equation}
 \boxed{\exp(tJ_\tau)=C_\tau(t)I+S_\tau(t)J_\tau.}
 \label{eq:euler-extendida-trit-fibonacci}
\end{equation}
In particular,
\[
 \exp(\pi J_{+1})+I=0,
 \qquad
 \exp(J_0)=I+J_0,
 \qquad
 \exp\!\bigl(2\log\varphi\,J_{-1}\bigr)=F_{\rm av}^{\,2}.
\]
The first equality is Euler's matrix closure; the second realizes the
parabolic threshold transport; the third incorporates Fibonacci
self-scaling through the eigenvalues \(\varphi\) and \(-\varphi^{-1}\).
\end{corollary}

\begin{proof}
The identity \eqref{eq:euler-extendida-trit-fibonacci} is obtained by separating
the even and odd powers of the exponential series and using
\(J_\tau^2=-\tau I\). The three specializations follow from
\(J_{+1}^2=-I\), \(J_0^2=0\), and from the already proved equality
\(F_{\rm av}^{\,2}=\exp(2\log\varphi\,J_{-1})\).
\end{proof}

\paragraph{\(3\) exact realizations of the quadratic regimes of the TRIT}
The elliptic, parabolic, and hyperbolic branches realize, respectively, the
internal Coxeter element of \(A_2\), the nilpotent generator, and the incidence
\(F_{\rm av}\) induced by the TPK calendar. The matrix identities
are exact in the declared spaces and gauges.

\paragraph{Realization conditions of the \(3\) quadratic regimes of TRIT and obstructions of algebraic nondegeneration}
The isolated TRIT type organizes the three quadratic possibilities, but it does not
by itself select the operators \(C\), \(N\), and \(F_{\rm av}\): the
selection uses, respectively, the APP Coxeter element, the nilpotent
transport model, and the TPK descent of modes. The appearance of \(\pi\) in the
elliptic spectral formula belongs to the analytic expression of the cycle and
does not replace the regional generation of \(\pi_{\rm HMT}\). The identity
\(A^{\mathsf T}GA=G\) constructs an exact Lorentz transformation in
dimension \(1+1\); it does not by itself identify this screen with
physical spacetime \(3+1\). A failure of
\[
(2C+I)^2=-3I,\qquad N^2=0,\qquad
(2A-3I)^2=5I,\qquad A^{\mathsf T}GA=G
\]
would refute the corresponding realization.
